% =====================================================================
% page-layout-scrlayer.tex — Running headers/footers wired via
% scrlayer-scrpage (KOMA-Script targets)
%
% Sibling of page-layout-fancyhdr.tex. Maps each of the twelve
% \keystone* content macros (defined by page-furniture.lua) to one of
% scrlayer-scrpage's parity-aware slot macros. Slot composition (text
% vs. page number, recto/verso fallback, center-slot conflict) lives in
% Lua. It also owns the head-rule default for this package: the rule is
% off by default, and the package-agnostic \keystoneheaderrule shim
% (defined below, emitted by page-furniture.lua only when the author
% sets header-rule: enabled) re-enables it through KOMA's headsepline.
% The fancyhdr counterpart defines the same shim over \headrulewidth.
%
% scrlayer-scrpage slot naming: the second character is page parity
% (o/e = odd/even = recto/verso); the third character is the column
% (l/c/r = left/center/right). The wirings below pair each parity
% slot with the matching \keystone*{L,C,R}{O,E} macro 1:1. Under
% scrartcl / scrreprt (single-sided by default), only the o-slots
% are selected; the e-slot wirings compile but never render.
%
% Chapter-start pages — the first page of every chapter, the title
% page, and the table of contents — render bare: no header, no
% footer, no rules. KOMA exposes \chapterpagestyle, \titlepagestyle,
% \partpagestyle, and \indexpagestyle as the macros that name which
% pagestyle the corresponding transitions select; we point each at
% the built-in "empty" pagestyle so those pages are bare. The plain
% pagestyle itself is not redefined here — KOMA's class-level
% transitions never select `plain` directly when these named macros
% are set, and authors who manually issue `\thispagestyle{plain}`
% in their content get scrlayer-scrpage's default plain pagestyle
% (which inherits from scrheadings).
% \@ifundefined guards each macro because scrartcl (no chapters) does
% not define \chapterpagestyle / \partpagestyle / \indexpagestyle.
%
% This file is part of the Keystone project layout.
% =====================================================================
\usepackage{scrlayer-scrpage}
\pagestyle{scrheadings}

% Head/foot rules off by default. The head rule is opt-in via the
% header-rule metadata key: page-furniture.lua emits \keystoneheaderrule
% — defined just below — only when the author enables it, flipping
% headsepline back on. (footsepline has no lever; both targets keep it
% off — see #391.)
\KOMAoptions{headsepline=false,footsepline=false}
\newcommand{\keystoneheaderrule}{\KOMAoptions{headsepline=true}}

% Dynamic-placeholder shims. The page-furniture filter emits these
% macro references when an author writes `{page}`, `{chapter}`, or
% `{section}` in `header-text` / `footer-text` content. KOMA's
% scrlayer-scrpage does not uppercase marks by default, so no
% `\nouppercase` wrapping is needed.
%
% `\leftmark` / `\rightmark` are explicit parity-independent mark
% accessors that return the chapter and section marks respectively,
% regardless of page parity. Under scrbook's default
% `\automark[chapter]{section}`, `\chapter` calls `\markboth{chapter}{}`
% (setting `\leftmark`) and `\section` calls `\markright{section}`
% (setting `\rightmark`). Mapping `{chapter}` → `\leftmark` and
% `{section}` → `\rightmark` keeps the semantics consistent with the
% fancyhdr companion — `{chapter}` is always the chapter title,
% `{section}` is always the section title, regardless of which side
% of a spread the page lands on. The context-sensitive `\headmark`
% would mix the two depending on page parity, which `\automark`'s
% interface exposes but isn't what authors writing `{chapter}` ask for.
\providecommand{\keystonepagemark}{\thepage}
\providecommand{\keystonechaptermark}{\leftmark}
\providecommand{\keystonesectionmark}{\rightmark}

\lohead{\keystoneheaderLO}\lehead{\keystoneheaderLE}
\cohead{\keystoneheaderCO}\cehead{\keystoneheaderCE}
\rohead{\keystoneheaderRO}\rehead{\keystoneheaderRE}
\lofoot{\keystonefooterLO}\lefoot{\keystonefooterLE}
\cofoot{\keystonefooterCO}\cefoot{\keystonefooterCE}
\rofoot{\keystonefooterRO}\refoot{\keystonefooterRE}

\makeatletter
\@ifundefined{chapterpagestyle}{}{\renewcommand*{\chapterpagestyle}{empty}}
\@ifundefined{titlepagestyle}{}{\renewcommand*{\titlepagestyle}{empty}}
\@ifundefined{partpagestyle}{}{\renewcommand*{\partpagestyle}{empty}}
\@ifundefined{indexpagestyle}{}{\renewcommand*{\indexpagestyle}{empty}}
\makeatother
